\section{The exact cost of initializing a stationary response}
\label{sec:r75-stationary-initialization}

The distinction between a freely chosen initial distribution and a
prescribed initial label is not peculiar to a periodic example.  There is
an exact general law.  Its proof is an elementary elimination in the
space of future response functions, within the classical positive
realization framework \cite{Heller65,MonrasWinter14,MonrasWinter16}.  It does not require
hidden states to lie in the minimal linear realization, or require their
unobserved transition matrix to be irreducible.

Let $A$ be a nonempty finite command alphabet, with a finite report set
$Y_a$ for each $a\in A$.  A finite classical realization consists of
nonnegative matrices $T_{a,y}$ such that
\[
 T_a:=\sum_{y\in Y_a}T_{a,y},\qquad T_a\one=\one,
\]
and an initial probability row $\mu$.  For a command--report word
$w=(a_1,y_1)\cdots(a_n,y_n)$ write
$T_w=T_{a_1,y_1}\cdots T_{a_n,y_n}$ and
$p(w)=\mu T_w\one$.  Optional terminal readouts have decoder
columns $d_{b,z}\ge0$, with $\sum_zd_{b,z}=\one$, and response
$p(w;b,z)=\mu T_wd_{b,z}$.  All following equalities include every specified
terminal readout when these are part of the experiment.

Call the prepared response \emph{predictively stationary under $\nu$}
if $\nu$ is a probability row on $A$ with every $\nu(a)>0$ and
\begin{equation}\label{eq:r75-predictive-stationarity}
 \sum_{a\in A}\nu(a)\sum_{y\in Y_a}p((a,y)w)=p(w)
 \qquad\text{for every future word }w.
\end{equation}
When terminal readouts are specified, additionally require
\begin{equation}\label{eq:r75-terminal-stationarity}
 \sum_{a\in A}\nu(a)\sum_{y\in Y_a}p((a,y)w;b,z)=p(w;b,z)
 \qquad\text{for all }w,b,z,
\end{equation}
including the empty word $w$.
This is an equality of observable responses; it does not assume
$\mu\sum_a\nu(a)T_a=\mu$ in a particular hidden realization.  For one
command and no terminal readout it is ordinary stationarity.  A response
has a one-label realization precisely when its report probabilities
factor as $p(w)=\prod_jp_{a_j}(y_j)$, and every terminal readout has a
fixed distribution independent of the preceding word.  For one command
this is the i.i.d. case.

If a stationary hidden row is desired, it can be chosen without adding
labels. Put $\overline T=\sum_a\nu(a)T_a$. Every $\mu\overline T^j$
has the same future response by the displayed stationarity identities.
A subsequential limit of the Ces\`aro averages
$H^{-1}\sum_{j=0}^{H-1}\mu\overline T^j$ is stationary for
$\overline T$ and retains that response. Thus the freely mixed minimum
below agrees with the convention requiring a stationary hidden row for
the averaged command, and with the usual stationary HMM convention when
there is only one command.

Let $n_{\mathrm{mix}}(p)$ be the minimum number of labels with an arbitrary
initial probability row, and let $n_{\mathrm{pure}}(p)$ be the minimum when
one prescribed initial label is counted.  The minima are $+\infty$ when
there is no finite realization.

\begin{theorem}[Stationary initialization law]
\label{thm:r75-stationary-initializer}
Suppose \eqref{eq:r75-predictive-stationarity} holds, together with
\eqref{eq:r75-terminal-stationarity} when terminal readouts are included.
If the response has
a one-label realization, both minima are one.  Otherwise
\begin{equation}\label{eq:r75-initialization-law}
 n_{\mathrm{pure}}(p)=n_{\mathrm{mix}}(p)+1,
 \qquad (+\infty)+1:=+\infty.
\end{equation}
The lower holds for all real nonnegative transition tables, including
reducible and nonminimal linear realizations.  The upper processes the
first actual command while entering the core; it neither discards a
physical call nor inserts an uncounted preliminary randomization.
\end{theorem}
\begin{proof}
Take any finite pure-start realization, with initial label $o$, and let
$f_i$ denote the complete future response from hidden label $i$.  Thus
$f_i(w)=e_i^{\mathsf T}T_w\one$, with the analogous values for terminal
readouts.  Put $\overline T=\sum_a\nu(a)T_a$.  Stationarity gives the
functional identity
\begin{equation}\label{eq:r75-initial-functional-identity}
 f_o=\sum_j\overline T_{oj}f_j.
\end{equation}
If $\overline T_{oo}=1$, nonnegativity and full support of $\nu$ imply
$(T_a)_{oo}=1$ for every $a$.  All report transitions from $o$ then return
to $o$.  Their probabilities, and the terminal readouts at $o$, constitute
a one-label realization.  In the case under consideration this is
excluded.  Hence
\[
 \beta_j=\frac{\overline T_{oj}}{1-\overline T_{oo}}\quad(j\ne o)
\]
is a probability row and $f_o=\sum_{j\ne o}\beta_jf_j$.

Delete label $o$.  On the remaining labels define
\begin{equation}\label{eq:r75-elimination-tables}
 \widehat T_{a,y}(i,j)
   =T_{a,y}(i,j)+T_{a,y}(i,o)\beta_j,
 \qquad i,j\ne o,
\end{equation}
and retain their terminal decoders.  These matrices are nonnegative and
their sums over $y$ are stochastic.  The recursion
$f_i((a,y)w)=\sum_jT_{a,y}(i,j)f_j(w)$, with the functional identity for
$f_o$ substituted on its right side, is exactly the recursion for
\eqref{eq:r75-elimination-tables}.  Induction on word length, starting with
the empty word and each terminal decoder, proves that every retained
$f_i$ is unchanged.  Initial row $\beta$ therefore realizes $f_o$ with
one fewer label.  This proves $n_{\mathrm{mix}}\le
n_{\mathrm{pure}}-1$ whenever a finite pure realization exists.

Conversely, given a $k$-label mixed-start realization $(\mu,T_{a,y})$,
append one label $o$.  Its report transition on command $a$ enters the
core with row $\mu T_{a,y}$ and never returns to $o$; the core retains its
original transitions.  A terminal readout at $o$ has row $\mu d_{b,z}$.
The first physical command is already included in this row, so induction
gives the original law at every word length, including zero.  This uses
$k+1$ labels and proves the reverse inequality.  A finite pure realization
is itself a finite mixed one, so the two finiteness assertions are
equivalent.  The one-label exception is immediate.
\end{proof}

\begin{corollary}[Stationary processes]
\label{cor:r75-stationary-process}
For every stationary finite-alphabet process that is not i.i.d., the
minimum pure-start edge-emitting hidden Markov realization has exactly
one more label than the minimum freely mixed realization, with the same
infinite-value convention.  In particular, the $q$ versus $q+1$ conclusion
of Theorem~\ref{thm:r74-depolarizing-obstruction} is an instance of a
general stationary principle, while its phase-dependent $q$ remains a
separate spectral conclusion.
\end{corollary}

Equality of these complete response functions also preserves the law
under every common randomized policy for commands, stopping and terminal
readout.  For each finite stopped history, multiply the equal response by
the same policy probabilities.  Summation gives all finite stopping
events, and decreasing continuation events give the nontermination
probability.  Consequently all nonnegative costs measurable from that
record have equal expectations, including infinite expectations.  These
are process-realization label counts.  If the external policy has its own
persistent memory, that memory is counted separately in a closed-loop
implementation; it is not supplied free by the theorem.
